% Owner: Prisca. Numbers: analysis/*.md (all recomputed from the raw run files).
\begin{abstract}
Large language models are increasingly used as code reviewers, and the code they review keeps changing. We ask
whether a change that leaves a vulnerability in place can make a model that correctly flagged it call the code
safe. We build 102 small JavaScript web handlers (17 vulnerable/safe pairs each for cross-site scripting, command
injection and path traversal) and an execution-based verifier that runs every program with attack inputs, so we
always know whether the bug is still there. On gpt-oss:20b, the only free model we tested that reliably recognizes
safe code, we apply four kinds of bug-preserving change. The model's own maintenance edits, nine
semantics-preserving ``testability tarpit'' transformations stacked up to eight deep, and larger multi-file
programs almost never fool it (0 of 162, 0 of 455 and 1 of 299 still-vulnerable programs called safe).
Deliberate attacks do: a comment that fakes the end of the code block gets a \safe{} verdict 62\% of the time, and
routing the input through an unseen helper with a sanitizer-like name 56\%. A delimited prompt reduces the first
attack to 3\% at no measurable cost, but leaves the second unchanged (54\%). A pilot on a commercial model
(gpt-6-luna) ignores the prompt injections entirely, yet still trusts the unseen helper 51\% of the time. Trusting
the name of code the model cannot see is the most robust failure we found.
\end{abstract}
